\clearpage
\fancyhead[L]{\small KNITLMAYER PLANETENABSTÄNDE.}
\hypertarget{native-533lower}{}\label{native-533lower}
\begin{center}\small Göttingische gelehrte Anzeigen. Stück 179. Seite 1777..1782. 1810 November 10.\\[3pt]\rule{.72\linewidth}{.3pt}\end{center}

\emph{Versuch einer genauen Darstellung des Progressionsverhältnisses der Planeten- und Trabantenabstände von ihren Centralkörpern.} Von \textsc{Ferdinand Knitlmayer}, Hauptmann in der kaiserl. königl. Oesterreichischen Armee. 1808. 48 Seiten in Octav. Brünn.

Der glückliche Erfolg, womit auf dem Wege der Erfahrung oder des Versuchens verschiedene Wahrheiten und Verhältnisse in der Astronomie gleichsam errathen worden sind, die sich nachher vollkommen bestätigt und zu wichtigen und grossen Folgerungen geführt haben, ist für viele Personen eine Aufmunterung gewesen, zwischen mancherlei in unserm Sonnensystem vorkommenden Zahlverhältnissen auf gut Glück Vergleichungen anzustellen und einen Zusammenhang zwischen ihnen zu suchen. Es gehören hieher hauptsächlich die mittlern Abstände der Planeten von der Sonne. In so fern die Planeten als einzelne schon gebildete Körper betrachtet werden, sind allerdings die Entfernungen derselben von dem Centralkörper sowohl, als alle übrigen Bestimmungsstücke ihrer Bahnen, als ganz von einander unabhängig und willkührlich anzusehen, und man hat durchaus keinen Grund, Relationen zwischen ihnen vorauszusetzen. Nimmt man indess an, dass unser Sonnensystem mechanischen Ursachen seine ursprüngliche Bildung verdanke, so liesse sich freilich die Möglichkeit denken, dass es zwischen Elementen der einzelnen Planeten und etwa ihren Massen und Dichtigkeiten Verhältnisse geben könnte, und falls man dergleichen wirklich entschieden wahrnähme, liesse sich vielleicht daraus mit Glück auf eine erste Entstehungsart unsers Planetensystems zurückschliessen, über welche wir eigentlich gar nichts haben, als vage, unzulässige Hypothesen. Inzwischen haben alle Bemühungen dieser Art, womit manche Liebhaber der Astronomie sich gequält haben, zu gar Nichts geführt, und schwerlich wird irgend ein mathematischer Astronom darin mehr, als eine verzeihliche, aber unfruchtbare Spielerei erkennen.

In Deutschland hat man eine Zeit lang viel Gewicht auf ein vermeintes Gesetz gelegt, in welches man die Planetenabstände zwingen wollte, indem man sie als Aggregate einer Constante und der Glieder einer geometrischen Progression betrachtete, die von Planet zu Planet auf das Doppelte steige. Inzwischen erhielt man dadurch keine Uebereinstimmung, sondern nur eine rohe Näherung; den Merkur musste man von der Reihe ganz absondern, da sein Abstand sich nicht jenem Aggregate, sondern der Constante selbst näherte, und für den Platz zwischen Mars und Jupiter haben sich sogar vier Prätendenten gefunden, von denen der am spätesten bekannt gewordene, der aber bei weitem der glänzendste, und wahrscheinlich der grösste von allen ist, die Vesta, gerade am beträchtlichsten von der Progression abweicht.

Die vorliegende kleine Schrift hat zur Absicht, ein anderes Gesetz aufzustellen, nach welchem die Planetenabstände sich nicht näherungsweise, sondern genau richten sollen. Hr. \textsc{Knitlmayer} geht von einer geometrischen Progression aus und nimmt an, dass die wirklichen mittlern Abstände der einzelnen Planeten sich von den Gliedern dieser Progression desto mehr entfernen, je grösser die Massen und je grösser der Unterschied der Dichtigkeit von der Dichtigkeit der Sonne sind, so dass die mittlern Abstände der verschiedenen Planeten durch die Formel
\SourceLeaf{534}
\[ a \,.\, 2^n + b\,m + c\,(D - \mathrm{1}) \]
dargestellt werden sollen, wo $m$ die Masse jedes Planeten; $D$ seine Dichtigkeit, die der Sonne als Einheit angenommen; $a$, $b$, $c$ constante Coëfficienten, und $n$ für Merkur 1, für Venus 2, für die Erde 3, für den Mars 4, für die neuen Planeten 5, für Jupiter 6, Saturn 7, Uranus 8 bedeuten.

Zur Prüfung einer solchen, an und für sich doch auf Nichts a priori gegründeten, sondern, um es mit dem rechten Worte zu bezeichnen, aus der Luft gegriffenen, Hypothese, sollte man denken, müssten die Werthe der Constanten $a$, $b$, $c$ aus den mittlern Entfernungen dreier Planeten, bei welchen Masse und Dichtigkeit am besten bekannt sind, also der Erde, des Jupiter und des Saturn, durch Elimination bestimmt, und dann die Formel an Planeten, wo Masse und Dichtigkeit leidlich genau bekannt sind, geprüft werden; also bei der Venus und dem Uranus; wobei man denn doch, weil wir von der Venus und vom Uranus die Dichtigkeit eigentlich nur obenhin kennen, immer noch sich hüten müsste, auf eine etwaige Uebereinstimmung zu viel Gewicht zu legen. Allein der Verfasser ist weit entfernt, auf eine solche Weise zu Werke zu gehen. Vielmehr nimmt er sofort $a = \mathrm{0{,}075}$ an, ohne dass dazu weiter ein Grund abgesehen werden kann, als weil man dieselbe geocentrische Progression bei der oben erwähnten Hypothese zum Grunde gelegt hatte. Die Coëfficienten $b$ und $c$ sind von ihm aus den mittlern Distanzen der Erde und des Jupiter bestimmt, und zwar, wenn wir die von ihm gebrauchte, grösstentheils sehr confuse, Darstellungsform in die unsrige übersetzen, so, dass $b = \mathrm{427{,}435}$, $c = \mathrm{0{,}135143}$ (die Masse der Sonne $= \mathrm{1}$ gesetzt). Bei den Planeten Venus, Mars, Saturn und Uranus nimmt er dann die Massen so an, wie \textsc{Wurm} und \textsc{Laplace} sie angegeben haben, und berechnet daraus nach seiner hypothetischen Formel die Dichtigkeiten, und daraus, in Verbindung mit den Massen, den körperlichen Inhalt und den scheinbaren Halbmesser in der Distanz 1. Dass diese, so wie sie vom Verf. berechnet sind, mit den beobachteten ziemlich nahe zusammenstimmen, sieht er als einen Beweis von der Richtigkeit seiner Hypothese an, und dann wird es ihm nicht schwer, bei dem Merkur und den vier neuen Planeten die an sich unbekannten Massen und Dichtigkeiten aus den von \textsc{Schröter} bestimmten scheinbaren Durchmessern abzuleiten. Beim ersten Anblick könnte jene angebliche Uebereinstimmung vielleicht für Manchen etwas Blendendes haben. Wir wollen daher nur die Art anführen, wie der Verf. bei seinen Rechnungen zu Werke geht. Bei der Vesta, dem Saturn und dem Uranus findet sich, nach der obigen Formel, für $D - \mathrm{1}$ ein negativer Werth $= -p$, daher der Verf. diese Planeten für weniger dicht erklärt, als die Sonne: allein anstatt nun die Dichtigkeit $= \mathrm{1} - p$ zu setzen, rechnet er so, als ob sie $= -\dfrac{1}{1+p}$ wäre. Nach dieser Probe von des Verf. mathematischen Kenntnissen wird man uns alle weitere Bemerkungen wohl erlassen.

Der Verf. wendet hierauf seine Hypothese auch auf die Trabantensysteme der Hauptplaneten an, und zwar, wenn wir nicht seinem eignen, zum Theil unverständlichen und widersprechenden, Vortrage folgen, sondern seine eigentliche Idee aus den wirklichen Anwendungen, die er davon gibt, aufklären, auf folgende Art. Er nimmt an, die mittlern Abstände der einzelnen zu Einem Hauptplaneten gehörenden Trabanten müssen gleichfalls durch die Formel
\[ \mathrm{0{,}075} \,.\, 2^n + \mathrm{427{,}435}\, m + \mathrm{0{,}135143}\,(D - \mathrm{1}) \]
dargestellt werden (wo $m$ die Masse, $D$ die Dichtigkeit, $n$ die Ordnungszahl jedes Trabanten bedeutet), nur müsse dann nicht, wie oben, die mittlere Distanz der Erde von der Sonne als Längeneinheit angenommen werden, sondern bei jedem Planeten eine andere, und zwar nach des Verf. Vorschrift $\dfrac{A\BareRoot[3]{M}}{B}$, wo $A$ den mittlern Abstand des Merkur, $B$ den des vom Trabanten begleiteten Hauptplaneten von der Sonne, $M$ des letztern Masse bedeutet (die Masse der Sonne immer $= \mathrm{1}$ gesetzt). Der Verf. scheint diesen Werth der Längeneinheit so gewählt zu haben, dass daraus eine leidliche Ueberstim\SourceLeaf{535}mung bei dem Trabanten der Erde erhalten wird. Allein da derselbe weiter keinen theoretischen Grund hat, und die Trabanten der andern Planeten, wo uns eine Kenntniss ihrer Dichtigkeit ganz abgeht, zu einer Prüfung der Hypothese nicht gebraucht werden können, so können wir in der ganzen Anwendung der Hypothese auf die Trabantensysteme nichts weiter, als ein leeres Spiel mit Zahlen, erkennen.
